% Generated by code/13_generate_exhibits.R
\begin{table}[H]
\centering
\footnotesize
\setlength{\tabcolsep}{3pt}
\caption[PCA of the country typology feature matrix]{PCA of the country typology feature matrix}
\label{tab:country_typology_pca}
\begin{tabular}{@{}lrrL{3.0cm}L{6.0cm}@{}}
\toprule
Component & Variance & Cumulative & Dominant dimension & Largest loadings \\
\midrule
PC1 & 46.9\% & 46.9\% & Fiscal scale & Behavioural effect; Preferred revenue; Revenue per resident; Scenario spread \\
PC2 & 23.4\% & 70.3\% & Legal allocation & Residual IIR/UTPR; QDMTT share; SbS extinguished share; Revenue/CIT \\
PC3 & 12.3\% & 82.6\% & Fiscal scale & Revenue/GDP; Revenue/CIT; Substance effect; High-risk imputation \\
\bottomrule
\end{tabular}

\vspace{0.35em}
\noindent\begin{minipage}{\textwidth}
\footnotesize\raggedright\textit{Notes:} The PCA is run on standardised country-level features. Revenue amounts, scenario spread and per-resident revenue enter in logs; signed main-effect amounts enter as signed log magnitudes. For PCA scaling, median imputation fills undefined channel shares for the four countries with zero assigned receipts.
\end{minipage}
\end{table}
